\subsection{Actual complete algebraic closures and rational lifted points}
\label{subsec:complete-algebraic-closures}

\begin{theorem}[Algebraic closedness from a genuinely dense subfield]
\label{spectrum:algebraic-closed-dense-completion}
\lean{EnvelopeProperty.isAlgClosed_of_denseRange_nontrivial,
EnvelopeProperty.isAlgClosed_of_denseRange,
EnvelopeProperty.UniformSpace.Completion.isAlgClosed_ultrametric}
\leanok
Let $K$ be an algebraically closed field, $L$ a complete ultrametric
normed field, and $K\to L$ an actual dense field homomorphism. Then $L$
is algebraically closed. Consequently the actual norm completion of any
ultrametric algebraically closed normed field is algebraically closed.
These statements impose neither characteristic zero nor nontriviality
of the norm.
\end{theorem}
\begin{proof}
First suppose the norm on $L$ is nontrivial. For a monic irreducible
$f\in L[T]$, choose an actual root $a$ in its finite splitting field,
endowed with the genuine spectral extension norm of the complete base.
The base field embeds isometrically into this normed splitting field,
so its image is closed by completeness.

Given any $\delta>0$, approximate the coefficients of $f$ by an actual
monic polynomial $g\in K[T]$ of the same positive degree, with error less
than
\[
 \frac{(\delta/\max\{|a|,1\})^{\deg f}}{\deg f+1}.
\]
The proved root-continuity estimate supplies a root of $g$ at distance
less than $\delta$ from $a$. Since $K$ is genuinely algebraically closed,
$g$ splits over $K$ and that root lies in the base-field image. Thus $a$
is in its closure and hence in the image itself. Pulling back the root
of $f$ proves algebraic closedness. This proof uses no separability or
Krasner lemma, so it also works in positive characteristic.

If the norm is trivial, its topology is discrete; density then makes the
actual coefficient homomorphism surjective. Its genuine field isomorphism
transfers roots. The completion corollary uses the actual dense canonical
completion homomorphism, the complete normed-field structure and the
proved preservation of ultrametricity.
\end{proof}

\begin{definition}[The actual complete algebraically closed valued extension]
\label{spectrum:complete-valued-algebraic-closure}
\lean{EnvelopeProperty.algebraicClosureValuation,
EnvelopeProperty.algebraicClosureValuation_algebraMap,
EnvelopeProperty.ValuedAlgebraicClosure,
EnvelopeProperty.valuedAlgebraicClosure_isAlgClosed,
EnvelopeProperty.valuedAlgebraicClosure_isUltrametricDist,
EnvelopeProperty.valuedAlgebraicClosure_norm_algebraMap,
EnvelopeProperty.valuedAlgebraicClosure_normedAlgebra,
EnvelopeProperty.CompleteAlgebraicClosure,
EnvelopeProperty.completeAlgebraicClosure_isAlgClosed,
EnvelopeProperty.completeAlgebraicClosure_norm_algebraMap}
\leanok
\uses{spectrum:complete-norm-extension,spectrum:algebraic-closed-dense-completion}
For every complete ultrametric normed field $K$, construct an actual
complete algebraically closed ultrametric normed field
\[
 F=\widehat{\overline K},
\]
whose actual coefficient embedding satisfies $|c|_F=|c|_K$ for all
$c\in K$. The construction includes trivial norms and arbitrary
characteristic and is independent of any analytic point space.
\end{definition}
\begin{proof}
The genuine extension theorem supplies a real-valued ultrametric
valuation on the actual algebraic closure of $K$ extending the entire
original norm. Choose that valuation and use its actual absolute value
to form the WithAbs normed field. Its genuine ring isomorphism to the
algebraic closure transfers algebraic closedness; the valuation inequality
proves ultrametricity. Its coefficient norm equality proves the actual
normed-algebra structure over $K$.

Take its actual uniform norm completion. Mathlib supplies the complete
normed-field structure and the inherited coefficient algebra. The preceding
algebraic-closed-completion theorem and ultrametric completion producer
give the claimed field properties. The canonical completion embedding
preserves norms, so the original coefficient embedding into $F$ still
preserves every original norm value. No extension field satisfying the
target properties is assumed as a certificate.
\end{proof}

\begin{theorem}[Actual rational rigid lift over the completed residue field]
\label{spectrum:rigid-complete-algebraic-closure-lift}
\leanok
\uses{spectrum:complete-valued-algebraic-closure,
spectrum:completed-residue-field,spectrum:relative-completed-residue-field,
spectrum:base-change-restriction}
For arbitrary commutative rings $R,A$, an $R$-algebra structure on $A$,
and an actual prescribed-base point $x$ of $A$, set
$F=\widehat{\overline{H(x)}}$ by the genuine construction above. There is
an actual $F$-rational point $x_F$ of $F\otimes_RA$, with the entire norm
of $F$ prescribed, which restricts to the whole original point $x$.
The residue quotient at its support is genuinely $F$; its support is
therefore maximal. No finite-type, coefficient-field or trivial-base
hypothesis is required.
\end{theorem}
\begin{proof}
The existing completed residue-field map sends $a\in A$ to an element
of $H(x)$ with norm $|a(x)|$. The actual compatible coefficient action
and tensor universal property give an $F$-algebra homomorphism
\[
 e_x:F\otimes_RA\longrightarrow F,
 \qquad c\otimes a\longmapsto c\,\iota(x(a)),
\]
where $x(a)$ denotes the actual completed-residue-field image and $\iota$
the constructed isometric field embedding. It is surjective, since
$e_x(c\otimes1)=c$.
Pull back the actual norm valuation of $F$ through this entire tensor
homomorphism. The coefficient inclusion is fixed, and on every pure
tensor the value is $|c|_F|a(x)|$. Thus the right restriction is exactly
the whole original point, including its prescribed coefficient values.
The norm of a field vanishes precisely at zero, so the support equals
$\ker e_x$. Surjectivity to the genuine field $F$ makes this kernel
maximal, and the actual first isomorphism theorem gives the asserted
whole residue-quotient $F$-algebra equivalence.

This constructs the complete algebraically closed extension and the
rational rigid point used in the choice of weights in Fang--Han--Xia
Proposition~5.22. The subsequent line, curve, reduction and normalization
arguments are separate producers and are not asserted here.
\end{proof}
